\documentclass[FGBook.tex]{subfiles}

\begin{document}

\proza{Zločin a trest}{FJODOR MICHALJEVIČ DOSTOJEVSKIJ}

\noindent\begin{wrapfigure}{r}{0.35\textwidth}\footnotesize
\subwection{Ukázka z díla:}
\setlength{\parindent}{3pt}
\textbfhl{Úder dolehl přímo na temeno}, její malý vzrůst tomu jen napomohl.
Vykřikla, ale docela slabounce, a znenadání se schoulila k~podlaze,
ačkoliv ještě stačila zdvihnout obě ruce.
V jedné pořád ještě svírala „zástavu“.
Tu \texthl{udeřil vší silou ještě jednou}, zase tupým koncem a zase do temene.
Krev vychlístla jako z převržené číše a tělo se~zvrátilo naznak.
Couvl před padajícím tělem a~ihned se sklonil k její tváři; \textbf{byla už mrtvá}.
Oči byly vytřeštěné, jako by~chtěly vylézt z důlků,
čelo a~celý obličej měla svraštělý a~zkřivený křečí.\\
Položil sekyru na zem vedle mrtvé a~rychle, ale~opatrně,
aby~se~nepotřísnil prýštící krví, strčil ruku stařeně do kapsy,
do~pravé kapsy, z~níž minule vytáhla klíče.
\textit{Byl naprosto příčetný}, mrákoty ani mdloby se~o~něj už nepokoušely,
jen \textit{ruce se mu pořád ještě třásly}.
Později se upamatoval, že byl dokonce velmi vnímavý
a~obezřetný a~že~neustále dával pozor, aby~se~nepotřísnil…
\end{wrapfigure}

\wection{LITERÁRNÍ TEORIE}

\subwection{Literární druh a žánr:}
\noindent Próza a psychologický román

% \subwection{Literární směr:}

% \subwection{Jazyk a slovní zásoba:}

\subwection{Figury:}
\noindent 
\enquote{Bože můj,} prosil, \enquote{ukaž mi cestu...} (\texthl{apostrofa}) \\
\enquote{a až zítra, zítra si všechno... Kdy jste přijely, už dávno?} (\texthl{epizeuxis a aposiopese})

\subwection{Tropy:}
\noindent 
Byl chudý jako kostelní myš (\texthl{přirovnání}) \\
Ach vy milá, nespravedlivá srdce! (\texthl{synekdocha})

\subwection{Stylistická charakteristika textu:}
\noindent 
Chronologický postup vyprávění děje er-formou od začátku do konce.
Výjimkou jsou \textit{retrospektivní sny a vzpomínky} u~kterých chronologický postup neplatí a \textit{vnitřní monology postav} u~kterých se přechází do ich-formy.
Text je rozdělen do \texthl{sedmi částí}, z nichž poslední je epilog.

\subwection{Postavy:}
\noindent 
\textschl{Rodion Romanovič Raskolnikov --} rozpolcená hlavní postava, jak už i trefné \textit{jméno napovídá}.
Jedná se o mladého muže, studenta práv, který žil dlouhou dobu v chudobě a celkově špatných podmínkách.
Začal tedy hledat určitý ideál a napadla ho \textbfhl{teorie prosazení zločinu}, kdy všichni géniové museli spáchat zločin, díky omezenosti zákonů a~prošlo jim to jen proto, že \textithl{dokázali, že mohou}.
Jeho nálada byla velice proměnlivá a často vedl monology. \\
\textschl{Dmitrij Prokofjič Razumichin --} skrz na skrz dobrý člověk.
Inteligentní, ale dobrý, díky tomu i trochu přihlouplý, protože \textit{odmítá pochopit myšlení vypočítavých lidí}.
Raskolnikovi pomáhá od začátku do konce a nakonec se ožení s Raskolnikovou sestrou, protože k ní během knihy zahoří láskou. \\
\textschl{Avdoťja Romanovna Raskolniková --} zkráceně Duňa, Raskolnikova sestra, která ho \texthl{bezmezně miluje} a pro zaručení zářné budoucnosti svému bratrovi, je ochotna vstoupit zcela dobrovolně i do manželství se zlým, však bohatým a vlivným mužem.
I když se jí nedostalo vzdělání jako Raskolnikovi, tak je hned v několika chvílích ukázáno, že je \texthl{neméně inteligentní}, než její bratr. \\
\textschl{Pulcherija Alexandrovna Raskolnikovová --} - matka Raskolnika, starší venkovská žena.
Svého syna miluje a to jí činí \textit{slepou k jeho problémům}.
I~když na konci je odhaleno, že přeci jenom \textbf{věděla více}, než kdy dala najevo. \\
\textschl{Sofja Semjonovna Marmeladová --} zkráceně Soňa, postava ze špatných poměrů, stejně jako Raskolnikov, kterou ale neopustila \texthl{pravá víra a~odevzdanost k Bohu} a~která nakonec Raskolnika \textbf{zachrání} právě díky Božímu učení. \\
\textschl{Arkadij Ivanovič Svidrigajlov --} stejně jako Soňa, i Svidrigajlov \textbfhl{ukazuje Raskolnikovi zrcadlo}.
Narozdíl od hodné Sofji, která stále věří v Boha a~neopouští jí to dobré, představuje Svidrigajlov co se z Raskolnika stane, když svůj \textit{hřích v sobě bude dusit}, či ho dokonce začne těšit.
Bezvěrec jenž má na svém krku nespočet životů a nakonec to ukončí sebevraždou.
Byl zamilován do Avdoťji. \\
\textschl{Porfirij Petrovič --} velice inteligentní vyšetřující soudce, který Raskolnika prokoukne hned zezačátku a po celou knihu mu činí \textbfhl{mentálního protivníka}.
Jejich dialogy jsou jedny z nejpovedenějších kousků.

\subwection{Děj:}
\noindent 
Příběh začíná u zdánlivě běžného dne v životě Raskolnika.
Pořád rozmyšlí nad něčím, co je mu natolik \textit{odporné}, že to neadresuje ani ve svých myšlenkách.
Uvažuje co ho k tomu vedlo.
Přitom \textbf{ochoří} a mysl se mu \textbf{zamlží}.
O několika dnech, při kterých trpěl horečkou, ani neví a nakonec spáchá~\enquote{to.}
Což~je~odhaleno, že se jedná o \textbfhl{vraždu} stařeny, která půjčuje peníze.
Ukradne peníze, ale schová je a jde se dál kurýrovat ze své nemoci.
Začíná \texthl{vyšetřování}, Raskolnikov je jeden z předních podezřelých.
Do života se mu vrací matka a sestra, přicházejí i noví lidé.
\textbfhl{Mezi nimi i Soňa}, která nakonec vyvede Raskolnika z temnoty, či \textbfhl{Svidrigajlov}, který po dlouhou dobu funguje spíš jako záhadná postava v pozadí.

\subwection{Kompozice, prostor a čas:}
\noindent 
První až šestá část se, s výjimkou vzpomínek a snových pasáží, odehrávají v Petrohradě, epilog pak na \textbfhl{Sibiři}.
Doba se dá \textit{odhadnout} na příbližně 2.~polovinu 19.~století.

\subwection{Význam sdělení:}
\noindent 
\textschl{Zločin a trest} je výstižným názvem pro celé téma.
Spáchá zločin a následuje trest.
Nejde mu uniknout, nejde jen o trest uložen soudem, ale~o~\texthl{trestání vlastním svědomím}, což je ukázáno hned na \texthl{několika postavách}, mezi nimi \texthl{hlavně Raskolnikov a Svidrigajlov}, kteří oba \texthl{zvolí rozdílnou cestu} ven.
Zatímco Svidrigajlov podléhá a únik již \textbfhl{vidí jen v sebevraždě}, Raskolnikov znovu \textbfhl{nalezne Boha a lásku k lidem} (převážně k Soně).

\wection{LITERÁRNÍ HISTORIE}

% \subwection{Politická situace:}

% \subwection{Základní principy fungování společnosti:}

% \subwection{Kontext dalších druhů umění:}

% \subwection{Kontext literárního vývoje:}

\subwection{Fjodor Michaljevič Dostojevskij {\ssmall -- život autora:}}
\noindent 
Ruský spisovatel, vystudoval vojenské technické učiliště.
S celou svojí skupinou byl \texthl{odsouzen na smrt} vojenským soudem, trest jim byl změněn na~nucené \texthl{práce na Sibiři}.
Prošel peklem a nápravu \textbf{shledal v Bohu}.
Ve svém okolí si všímal jen toho špatného, čehož také využíval ve svých dílech.
Věřil, že spása se dá najít jedině v Bohu a to také hlásal dál.

% \subwection{Vlivy na dané dílo:}

\subwection{Další autorova tvorba:}
\noindent 
Mezi nejznámnější díla patří \textsc{Zápisky z mrtvého domu}, právě \textsc{Zločin a trest}, \textsc{Idiot Idiot}, či \textsc{Bratři Karamazovi}

% \subwection{Jak dílo inspirovalo další vývoj literatury:}

% \wection{LITERÁRNÍ KRITIKA}

% \subwection{Dobové vnímání díla a jeho proměny:}

% \subwection{Aktuálnost tématu a zpracování tématu:}


\wection{OSOBNÍ NÁZOR}
\noindent 
Kniha je napsaná poutavě, postavy, které ani ve svých myšlenkách neadresují jádro myšlenky mají pro takové chování své důvody a nikoliv jen pro napětí čtenáře.
Myšlenka je tu s námi věčně a Dostojevskij ji dokonale \textbf{jmenuje}.

\vfill

\noindent\begin{minipage}{\textwidth}
    {\textcolor{\wpagecolor}{\rule{\linewidth}{0.4pt}}
    \footnotesize
    \textbfhl{Apostrofa --} řečnická figura, ve které postava mluví k nepřítomné osobě.
    V tomto případě k Bohu.
    \textbfhl{Epizeuxis --} básnická figura, opakování stejného slova či~slovního spojení. 
    \textbfhl{Aposiopese --} literární pojem označující odmlku, nedokončení věty.
    \textbfhl{Synekdocha --} rétorická figura při níž se celek označí názvem jeho části, či část názvem celku.
    Spřízněná s metonymií.
    \textbfhl{Metonymie --} přenos označení.
    Nazvání objektu názvem jiného objektu, dle určité souvislosti.
    }
\end{minipage}

\newpage

\end{document}